\chapter{背景诊断与分通道投影}
\section{空白段回归}
若目标通道受强背景物种影响，可在事先指定的空白窗口拟合
\begin{equation}
 y_q(t)=\beta_0+\sum_{r\in\mathcal R}\beta_ry_r(t)
 +\beta_pp(t)+\beta_tt+e_q(t).
\end{equation}
参考通道集合 $\mathcal R$、是否使用压力、是否使用线性漂移，都应由实验语境指定。预测量只在有效空白样点上计算中心和尺度，拟合后再向整个窗口预测。这样避免把反应段信息偷偷用于定义背景。

必须区分统计关系和物理校正。压力变量帮助解释共同波动，并不意味着已经获得了压力对检测灵敏度的传递函数。若压力与产物源项共同变化，回归也可能消去真实信号。软件因此保留模型与窗口信息，并输出未裁剪的正、负残差。
\fig{08_background_raw}{.91}{44、45、46 通道的合成背景记录。为比较相对变化，图中各曲线除以初始空白均值；回归仍使用原始数值。45 与 46 含有参考通道贡献、压力项和附加残差，不表示真实产物峰。}
\fig{09_residual_45}{.91}{质量 45 的合成带符号残差。压力辅助模型在此构造案例中更接近注入残差；这不是“压力校正后即为产物”的一般规则。}
\fig{09_residual_46}{.91}{质量 46 的合成带符号残差。保留负残差有助于检查背景失配及信号减小，程序不因非负源先验而修改这些观测。}

\section{联合求解中的背景变量}
另一个问题是，在反卷积中同时估计有限维背景。给定公共时间基 $B\in\R^{N\times p}$，每个通道有独立系数 $\beta_m$，考虑
\begin{equation}
\min_{x,\{\beta_m\}}\frac12\sum_m
\|W_m[(Ax)_m+B\beta_m-y_m]\|_2^2+\mathcal R(x).\label{eq:jointbg}
\end{equation}
本版提供常数基或常数加线性漂移基。其自由度有限，不能任意追踪反应峰。对每通道的加权基作薄 SVD：
\begin{equation}
 W_mB=Q_m\Sigma_mV_m^T,
\end{equation}
只保留通过秩阈值的列。$Q_m$ 的列正交，定义 $P_m=I-Q_mQ_m^T$。

\begin{proposition}[Gaussian 背景消元]
对固定 $x$，式~\eqref{eq:jointbg} 关于 $\beta_m$ 的最小值等于
\begin{equation}
\frac12\sum_m\|P_mW_m[(Ax)_m-y_m]\|_2^2+\mathcal R(x).
\end{equation}
一个最小范数系数解为
\begin{equation}
\beta_m=V_m\Sigma_m^{-1}Q_m^TW_m[y_m-(Ax)_m].\label{eq:betarecover}
\end{equation}
\end{proposition}
\begin{proof}
把 $W_m[y_m-(Ax)_m]$ 分解为 $Q_m$ 列空间上的正交投影和正交补分量。前者可由 $W_mB\beta_m$ 完全表示，后者与该列空间正交，无法由任何 $\beta_m$ 改变。勾股分解给出最小残差；对保留的非零奇异值求逆即得式~\eqref{eq:betarecover}。
\end{proof}
这属于变量投影思想在线性背景参数上的直接应用\cite{varpro}。本项目的 PARS 命名指投影、加速和正则化求解的工程组合，并不是宣称发现新的变量投影原理。

\section{投影后的算子与严格伴随}
令 $T=PW A$，其中 $P$ 和 $W$ 按通道作用。因 $P_m^T=P_m$，有
\begin{equation}
T^*=A^*WP.
\end{equation}
注意一般 $W_mP_m\ne P_mW_m$；实现中不能任意交换两个步骤。新求解器以 $T$ 和 $PWy$ 调用 FISTA，结束后恢复背景系数，再用原始前向模型重建信号和带符号残差。

\begin{algorithm}[htbp]
\caption{分通道背景消元的谱--时间求解}
\KwIn{$Y,S,h_j,B,W_m$，正则化系数、停止阈值与迭代上限}
\KwOut{源项 $X$、背景、重构、残差与诊断}
对各通道计算 $W_mB$ 的秩截断 SVD，保存 $Q_m,\Sigma_m,V_m$\;
建立批量谱--时间算子 $A$，定义 $T=PW A$ 及 $T^*=A^*WP$\;
检查伴随内积恒等式，计算投影目标 $PWY$\;
用回溯与重启保护的 FISTA 最小化投影目标加正则项\;
按式~\eqref{eq:betarecover} 恢复每通道背景系数\;
计算原始域重构 $AX+B\beta$、残差和投影正交性\;
分别保存任务完成状态与数值收敛状态，不因到达上限而改写为收敛\;
\end{algorithm}

\section{为什么按通道构造}
直接把所有背景拼为整体块对角矩阵，其形状为 $(MN)\times(Mp)$；若以稠密数组存储，需要 $M^2Np$ 个浮点数。实际上各通道彼此独立，只需保存 $M$ 个 $N\times p$ 因子，主存储量为 $MNp$。

对于本次 $M=24,N=2048,p=2$ 的基准，结构化因子存储约为稠密实现的 $1/23.96$，暖投影调用约快 5.57 倍。这里比较的是同一投影运算与因子存储，不是整个程序的峰值内存。SVD 构造与求解器运行仍有其它开销。

\section{不可辨识方向与适用限制}
若存在非零 $v$ 使 $WAv$ 落入背景子空间，则 $PWAv=0$。此方向上的源变化可以被背景系数抵消，数据无法单独确定它。正则化只能从这些等价解中偏好某一解，不能把偏好解释为实验识别。

例如，一个持续不变的产物与常数背景在稳态区间可能混淆；传输核已知也不必然消除该混淆。应使用可解释的空白、切换实验、独立标定或限制背景自由度来增加证据。本报告的模拟特意包含源项为零的边界区间，有利于区分背景与脉冲，但这不是所有实验都具备的条件。

上述解析消元依赖平方损失。Huber 目标关于背景不再由同一个固定正交投影精确消元，因此本版 \code{channel_pars} 明确拒绝 Huber 参数，不把近似结果冒充上述命题的精确实现。
